% ╔══════════════════════════════════════════════════════════════════════════╗
% ║   ÍNDICE DETALLADO — ECONOMÍA PREUNIVERSITARIA                           ║
% ║   8 Unidades · 56 Temas                                                  ║
% ╚══════════════════════════════════════════════════════════════════════════╝

{
\hypersetup{linkcolor=AzulPizarra}
\tableofcontents
}
\newpage

\thispagestyle{fancy}

\begin{center}
  {\Large\sffamily\bfseries\color{AzulPizarra} Índice Temático — Economía Preuniversitaria}\\[4pt]
  {\small\sffamily\color{GrisMedio}
  8 Unidades · 56 Temas · Micro, Macro, Finanzas, Economía Peruana e Historia del Pensamiento}
\end{center}

\vspace{0.5cm}

% ── UNIDAD I ──────────────────────────────────────────────────────────────
\begin{tcolorbox}[enhanced, colback=AzulPizarra, colframe=DoradoAcadémico,
    arc=3pt, boxrule=1pt, left=8pt, right=8pt, top=5pt, bottom=5pt,
    fontupper=\sffamily\bfseries\color{white}]
  UNIDAD I \hfill Introducción a la Economía \hfill Temas 1–7
\end{tcolorbox}
\begin{multicols}{2}
  \begin{enumerate}[label=\textcolor{AzulPizarra}{\textbf{\arabic*.}},
      leftmargin=2em, itemsep=1pt]
    \item Concepto de economía
    \item Necesidades humanas
    \item Bienes y servicios
    \item Escasez
    \item Costo de oportunidad
    \item Agentes económicos
    \item Actividades económicas
  \end{enumerate}
\end{multicols}

\vspace{0.3cm}

% ── UNIDAD II ─────────────────────────────────────────────────────────────
\begin{tcolorbox}[enhanced, colback=AzulMedio, colframe=DoradoAcadémico,
    arc=3pt, boxrule=1pt, left=8pt, right=8pt, top=5pt, bottom=5pt,
    fontupper=\sffamily\bfseries\color{white}]
  UNIDAD II \hfill Sistemas Económicos \hfill Temas 8–11
\end{tcolorbox}
\begin{multicols}{2}
  \begin{enumerate}[label=\textcolor{AzulMedio}{\textbf{\arabic*.}},
      leftmargin=2em, itemsep=1pt, start=8]
    \item Sistema capitalista
    \item Sistema socialista
    \item Sistema mixto
    \item Doctrinas económicas
  \end{enumerate}
\end{multicols}

\vspace{0.3cm}

% ── UNIDAD III ────────────────────────────────────────────────────────────
\begin{tcolorbox}[enhanced, colback=AzulPizarra, colframe=DoradoAcadémico,
    arc=3pt, boxrule=1pt, left=8pt, right=8pt, top=5pt, bottom=5pt,
    fontupper=\sffamily\bfseries\color{white}]
  UNIDAD III \hfill Microeconomía \hfill Temas 12–20
\end{tcolorbox}
\begin{multicols}{2}
  \begin{enumerate}[label=\textcolor{AzulPizarra}{\textbf{\arabic*.}},
      leftmargin=2em, itemsep=1pt, start=12]
    \item Mercado
    \item Oferta
    \item Demanda
    \item Elasticidad
    \item Equilibrio de mercado
    \item Competencia perfecta
    \item Monopolio
    \item Oligopolio
    \item Competencia monopolística
  \end{enumerate}
\end{multicols}

\vspace{0.3cm}

% ── UNIDAD IV ────────────────────────────────────────────────────────────
\begin{tcolorbox}[enhanced, colback=AzulMedio, colframe=DoradoAcadémico,
    arc=3pt, boxrule=1pt, left=8pt, right=8pt, top=5pt, bottom=5pt,
    fontupper=\sffamily\bfseries\color{white}]
  UNIDAD IV \hfill Macroeconomía \hfill Temas 21–28
\end{tcolorbox}
\begin{multicols}{2}
  \begin{enumerate}[label=\textcolor{AzulMedio}{\textbf{\arabic*.}},
      leftmargin=2em, itemsep=1pt, start=21]
    \item Producto Bruto Interno
    \item Inflación
    \item Deflación
    \item Recesión
    \item Desempleo
    \item Ciclos económicos
    \item Crecimiento económico
    \item Desarrollo económico
  \end{enumerate}
\end{multicols}

\vspace{0.3cm}

% ── UNIDAD V ──────────────────────────────────────────────────────────────
\begin{tcolorbox}[enhanced, colback=AzulPizarra, colframe=DoradoAcadémico,
    arc=3pt, boxrule=1pt, left=8pt, right=8pt, top=5pt, bottom=5pt,
    fontupper=\sffamily\bfseries\color{white}]
  UNIDAD V \hfill Sistema Financiero \hfill Temas 29–34
\end{tcolorbox}
\begin{multicols}{2}
  \begin{enumerate}[label=\textcolor{AzulPizarra}{\textbf{\arabic*.}},
      leftmargin=2em, itemsep=1pt, start=29]
    \item Dinero
    \item Sistema financiero
    \item Bancos
    \item Crédito
    \item Tasas de interés
    \item Bolsa de valores
  \end{enumerate}
\end{multicols}

\vspace{0.3cm}

% ── UNIDAD VI ────────────────────────────────────────────────────────────
\begin{tcolorbox}[enhanced, colback=AzulMedio, colframe=DoradoAcadémico,
    arc=3pt, boxrule=1pt, left=8pt, right=8pt, top=5pt, bottom=5pt,
    fontupper=\sffamily\bfseries\color{white}]
  UNIDAD VI \hfill Economía Peruana \hfill Temas 35–43
\end{tcolorbox}
\begin{multicols}{2}
  \begin{enumerate}[label=\textcolor{AzulMedio}{\textbf{\arabic*.}},
      leftmargin=2em, itemsep=1pt, start=35]
    \item BCRP
    \item SUNAT
    \item MEF
    \item INEI
    \item Política fiscal
    \item Política monetaria
    \item Comercio exterior
    \item Globalización
    \item TLC
  \end{enumerate}
\end{multicols}

\vspace{0.3cm}

% ── UNIDAD VII ────────────────────────────────────────────────────────────
\begin{tcolorbox}[enhanced, colback=AzulPizarra, colframe=DoradoAcadémico,
    arc=3pt, boxrule=1pt, left=8pt, right=8pt, top=5pt, bottom=5pt,
    fontupper=\sffamily\bfseries\color{white}]
  UNIDAD VII \hfill Historia del Pensamiento Económico \hfill Temas 44–49
\end{tcolorbox}
\begin{multicols}{2}
  \begin{enumerate}[label=\textcolor{AzulPizarra}{\textbf{\arabic*.}},
      leftmargin=2em, itemsep=1pt, start=44]
    \item Mercantilismo
    \item Fisiocracia
    \item Escuela clásica
    \item Marxismo
    \item Keynesianismo
    \item Neoliberalismo
  \end{enumerate}
\end{multicols}

\vspace{0.3cm}

% ── UNIDAD VIII ───────────────────────────────────────────────────────────
\begin{tcolorbox}[enhanced, colback=AzulMedio, colframe=DoradoAcadémico,
    arc=3pt, boxrule=1pt, left=8pt, right=8pt, top=5pt, bottom=5pt,
    fontupper=\sffamily\bfseries\color{white}]
  UNIDAD VIII \hfill Actualidad Económica \hfill Temas 50–56
\end{tcolorbox}
\begin{multicols}{2}
  \begin{enumerate}[label=\textcolor{AzulMedio}{\textbf{\arabic*.}},
      leftmargin=2em, itemsep=1pt, start=50]
    \item Pobreza
    \item Desigualdad
    \item Desarrollo sostenible
    \item Economía ambiental
    \item Crisis económicas
    \item Economía digital
    \item Inteligencia artificial y economía
  \end{enumerate}
\end{multicols}

\vspace{0.5cm}

\begin{cajatip}
  En los exámenes de admisión de ciencias sociales, los temas de mayor frecuencia
  son: Oferta y Demanda (Unidad III), PBI e Inflación (Unidad IV), y Economía
  Peruana, especialmente BCRP y SUNAT (Unidad VI). Se recomienda manejar con
  fluidez los cuadros comparativos de sistemas económicos y estructuras de mercado.
\end{cajatip}

\newpage